\section{Discussion}
\label{sec:discussion}

The result (\S\ref{sec:results}) is neither cleanly ``portable'' nor
cleanly ``fragile'' -- it is source-pair-dependent in a specific,
mechanistically legible way, and that dependence itself is the finding.

\subsection{Portability Is Real but Source-Pair-Dependent, Not Uniform}

Comparative scheduler rankings are never uncorrelated between any two
sources at any region ($\tau \geq 0.55$ everywhere,
\S\ref{sec:cross-source-rankings}): a scheduler comparison performed on
one of this study's three sources is never scientifically meaningless
evidence about the others. At the same time, the degree of portability is
far from uniform, and the source of the non-uniformity is legible rather
than diffuse: Azure-2024 and Bailian/Qwen agree with each other almost
perfectly at every operating region ($\tau = 0.89$--$1.00$), while every
comparison involving BurstGPT sits $0.2$--$0.4$ lower in Kendall's tau-b,
consistently across all six regions (\S\ref{sec:cross-source-rankings},
\S\ref{sec:results-robustness}'s leave-one-source-out breakdown). The
practical implication is neither ``a small, well-chosen evaluation
footprint always suffices'' nor ``every scheduler comparison must be
re-verified on every workload'': it is that \emph{which} second source a
comparison is checked against matters more than \emph{whether} a second
source is checked at all. A comparison validated only on Azure-2024-like
traffic transfers with high confidence to Bailian/Qwen-like traffic under
this study's panel and metric, but the same transfer to BurstGPT-like
traffic is measurably less reliable, and that reliability gap holds at
every load level tested, not only under stress.

\subsection{Practically Supported Reversals Are Real, Uncommon, and
Load-Concentrated}

Of 990 pairwise (policy pair $\times$ condition) comparisons on the
PRIMARY metric, $3.6\%$ are practically and statistically supported
reversals (\S\ref{sec:reversals}) -- a small minority, but not zero, and
therefore not dismissible as measurement noise: the same bootstrap
procedure that supports these 36 reversals also classified 906 comparisons
as stable and produced zero comparisons in the ambiguous
\texttt{UNSUPPORTED\_SIGN\_CHANGE\_WIDE\_CI} class. Every supported
reversal we examined is the same policy pair
(\texttt{slai\_faithful} vs.\ \texttt{vllm\_faithful}) and involves
BurstGPT on one side, and reversal counts rise from zero at \texttt{LOW}
to a peak at \texttt{OVERLOAD} (\S\ref{sec:reversals}) -- reversals are not
spread evenly across the operating envelope; they concentrate at and above
\texttt{PRE\_KNEE}. A scheduler comparison performed only at low load,
even on a source pair that is otherwise highly portable, would have missed
every reversal this benchmark found. This motivates evaluating scheduler
comparisons across an operating curve, not at a single convenient load
point, as this benchmark's design already does
(\S\ref{sec:calibration}).

\subsection{Practical Implications for Benchmark Design}

Two implications follow directly and are not deployment recommendations
about which scheduler to use, only about how to evaluate schedulers
credibly. First, a comparison checked against only one additional workload
source cannot be assumed representative of ``other workloads'' in
general -- which specific source is used as the check matters, and
BurstGPT-like traffic in particular is where this study's evidence
concentrates its cross-source disagreement. Second, benchmark sample
complexity is itself load- and metric-behavior-relevant evidence:
\S\ref{sec:sample-complexity}'s finding that exact-order recovery requires
the full window budget for two of three sources, while top-1 recovery is
achievable from very few windows for those same two sources, means a
researcher who only needs to know ``which scheduler is best'' faces a
different, much smaller evidence requirement than one who needs the
complete ordering -- a distinction a single point-estimate benchmark
result cannot convey on its own.

\subsection{Explanatory Telemetry: A Preliminary, Heavily Qualified Read}

The frozen, pre-specified explanatory model (\S\ref{sec:telemetry},
\texttt{telemetry\_explanation.json}) associates the five workload
descriptors with reversal sites at each (policy pair, condition) cell, but
converges with a usable sample size ($\geq 20$ observations) for only 42 of
990 cells -- most reversal sites are too sparse for this per-cell model to
fit reliably, and the frozen output records point coefficients only, with
no standard errors or significance test. Any reading of this evidence must
therefore be provisional. Within that small, non-significance-tested set,
\texttt{prompt\_tokens\_cv} has a negative coefficient in 39 of 42 cells --
the most consistent single direction among the five descriptors -- which
is at least \emph{consistent with} (not confirmatory of) prompt-length
variability being associated with where reversals occur, echoing
\S\ref{sec:drivers}'s finding that prompt-length structure is the primary
driver of source separability itself. The other four descriptors show
weaker, closer-to-even sign splits across the interpretable cells and do
not support a directional claim. We do not read this as evidence that any
descriptor \emph{causes} a reversal, only as a candidate association worth
a properly powered, purpose-built follow-up study; consistent with
\S\ref{sec:evidence-boundaries}, this model is explanatory only and is
never used to select, route, or recommend a scheduler.

\subsection{Workload Descriptor Interpretation}

The workload-separability result (\S\ref{sec:separability}) establishes
that the three sources are observably different, and that prompt-length
structure is the primary driver of that separability
(\S\ref{sec:drivers}). Whatever RQ4's descriptor-to-reversal association
analysis finds, it should be read against this backdrop: an association
between a descriptor and a reversal does not by itself establish that the
descriptor is a genuine \emph{causal} driver of the reversal, only that it
is a source-separating feature that co-occurs with where reversals were or
were not observed.

\subsection{Benchmark Construction and Practical Reproducibility}

Regardless of outcome direction, this study's paired design, fixed
scheduler panel, and preregistered protocol (\S\ref{sec:study-design})
are intended to make the benchmark itself reusable: a future scheduler
design can be dropped into the same panel and evaluated against the same
frozen windows, operating regions, and metric contract, using the planned
public artifact (\S\ref{sec:artifact}) rather than re-deriving the protocol
from this manuscript's prose.
